% Generated draft: requires mathematical and source review.
\documentclass[11pt]{article}
\usepackage[margin=25mm]{geometry}
\usepackage{fontspec}
\setmainfont{Latin Modern Roman}
\usepackage{amsmath,amssymb,mathtools,unicode-math}
\setmathfont{Latin Modern Math}
\usepackage{enumitem,needspace,xurl,hyperref}
\setlength{\parindent}{0pt}
\setlength{\parskip}{5pt}
\setlength{\emergencystretch}{3em}
\newcommand{\N}{\mathbb{N}}
\newcommand{\Z}{\mathbb{Z}}
\newcommand{\Q}{\mathbb{Q}}
\newcommand{\R}{\mathbb{R}}
\newcommand{\C}{\mathbb{C}}
\newcommand{\F}{\mathbb{F}}
\newcommand{\A}{\mathbb{A}}
\newcommand{\PP}{\mathbb{P}}
\newcommand{\E}{\mathbb{E}}
\newcommand{\eps}{\varepsilon}
\newcommand{\dd}{\,\mathrm{d}}
\begin{document}
\section*{UnsolvedMath ID: 20001284}
\textbf{Source problem code:} catalogue entry
\subsection{Definitions and mathematical statement}
An origin-symmetric star body \(K\subset\R^3\) is described by a positive continuous even radial function \(\rho_K:S^2\to(0,\infty)\), with \(K=\{tu:0\le t\le\rho_K(u),\ u\in S^2\}\). For a plane \(u^\perp\) through the origin, let \(P_K(u)\) be the Euclidean length of the boundary of \(K\cap u^\perp\), when this boundary is rectifiable. Gardner's question asks whether
\[
P_K(u)=P_L(u)\quad\text{for every }u\in S^2
\quad\Longrightarrow\quad K=L
\]
for origin-symmetric star bodies \(K,L\) in \(\R^3\). The source does not impose regularity or explicitly require finite section perimeters; those qualifications require checking for the full stated class, since continuous radial functions alone do not guarantee rectifiable section boundaries. No convexity assumption is silently added.

The supplied verification identifies this as Problem 20, not the extracted ``0,'' of the 2010 workshop list. The source cites [G, Problem 7.6] and reports affirmative results for bodies of revolution by R. Howard, F. Nazarov, D. Ryabogin and A. Zvavitch and for polytopes by V. Yaskin. These are source-reported special cases, not a global resolution or independently verified current literature assessment.

The research aid gives a finite-dimensional local theorem. Let \(m\ge0\) be an integer, let \(\mathcal H_\ell\) be the real spherical harmonics of degree \(\ell\) on \(S^2\), and put
\[
V_m=\bigoplus_{j=0}^m\mathcal H_{2j},\qquad D=\dim V_m=(m+1)(2m+1).
\]
For every \(R>0\), consider the even smooth radial chart \(\rho=R+f\), \(f\in V_m\), sufficiently small that \(\rho>0\). If \(C_u=S^2\cap u^\perp\) and \(\partial_\tau\) is differentiation along its unit-speed circular parameter, its section-perimeter measurement is
\[
P_{R+f}(u)=\int_{C_u}\sqrt{(R+f)^2+(\partial_\tau f)^2}\,\dd s.
\]
There exist \(D\) plane normals \(u_1,\ldots,u_D\) such that the nonlinear measurement map
\[
f\longmapsto\bigl(P_{R+f}(u_1),\ldots,P_{R+f}(u_D)\bigr)
\]
is locally invertible near \(f=0\) as a map between real vector spaces of dimension \(D\). In particular, for some neighborhood \(U\subset V_m\) of zero, equality of these \(D\) measurements for \(f,g\in U\) implies \(f=g\). This is local uniqueness within the fixed even bandlimited model; the neighborhood and model restriction are essential.

The aid identifies the first derivative at the ball as the Funk transform and the second derivative as
\[
\mathrm DP_R[f](u)=\int_{C_u}f\,\dd s,\qquad
\mathrm D^2P_R[f,g](u)=\frac1R\int_{C_u}(\partial_\tau f)(\partial_\tau g)\,\dd s.
\]
This finite-dimensional result does not prove uniqueness for arbitrary star bodies, for unrestricted smooth perturbations, or away from a neighborhood of a ball.

The registry labels the entry partially solved; attempt 1 is a partial result, with open status at the run, internal rigor label proved and novelty confidence low. These labels accompany a machine-generated research aid that is not peer reviewed and requires independent verification before citation.
\subsection{Short English statement}
Gardner asks whether the perimeters of all central planar sections determine an origin-symmetric star body in three dimensions. The supplied aid proves a local result near a ball when the radial function lies in a fixed finite spherical-harmonic band: finitely many suitably chosen plane measurements determine the body within that model. The unrestricted question and its perimeter regularity requirements remain separate.
\subsection{Sources}
\begin{itemize}
\item Local UnsolvedMath registry, permanent ID \(20001284\), source code \texttt{catalogue entry}; \texttt{aim-convex-geometry-notes.json}, \texttt{notes[15]}.
\item AIM, \emph{Problems from the workshop ``Mahler's conjecture and duality in convex geometry''}, August 9--13, 2010, notes by Jaegil Kim, recovered Problem 20 (R. J. Gardner): \url{https://aimath.org/WWN/mahlerduality/mahlerduality.pdf}. The numbering recovery is attributed to the supplied verification, not independent source inspection.
\item [G, Problem 7.6], as printed. The record does not supply the bibliographic identity of [G]. It attributes the bodies-of-revolution case to R. Howard, F. Nazarov, D. Ryabogin and A. Zvavitch, and the polytopal case to V. Yaskin, without full references.
\item Supplied research attempt 1 and registry update August 14, 2026; no separate authoritative citation for the finite-dimensional local theorem is provided. The original literature field is empty.
\end{itemize}
\end{document}
